\chapter{Módulo 7 --- Estadística y modelos en R}

\begin{itemize}
\tightlist
\item
  \textbf{Duración estimada:} 1.5 h teoría + 2.5 h práctica
\item
  \textbf{Objetivos de aprendizaje.} Al finalizar, la persona podrá:

  \begin{enumerate}
  \def\labelenumi{\arabic{enumi}.}
  \tightlist
  \item
    \textbf{Calcular} estadísticos descriptivos (media, mediana,
    desviación estándar, cuantiles, correlación).
  \item
    \textbf{Ejecutar} pruebas de hipótesis básicas
    (\passthrough{\lstinline!t.test()!},
    \passthrough{\lstinline!wilcox.test()!},
    \passthrough{\lstinline!chisq.test()!}).
  \item
    \textbf{Ajustar} un modelo de regresión lineal con
    \passthrough{\lstinline!lm()!} usando la sintaxis de fórmulas.
  \item
    \textbf{Interpretar} coeficientes, valores p, R² y residuos, y
    \textbf{justificar} si el modelo es adecuado.
  \end{enumerate}
\end{itemize}

\section{Contenidos}

\begin{enumerate}
\def\labelenumi{\arabic{enumi}.}
\tightlist
\item
  Estadística descriptiva por grupo con \passthrough{\lstinline!dplyr!}.
\item
  Correlación: \passthrough{\lstinline!cor()!} (Pearson y Spearman).
\item
  Pruebas de hipótesis: t de Welch, Mann-Whitney/Wilcoxon, chi-cuadrada.
\item
  Fórmulas en R: \passthrough{\lstinline!y \~ x1 + x2!}.
\item
  Regresión lineal: \passthrough{\lstinline!lm()!},
  \passthrough{\lstinline!summary()!}, \passthrough{\lstinline!coef()!},
  \passthrough{\lstinline!confint()!},
  \passthrough{\lstinline!predict()!}.
\item
  Diagnóstico con residuos; resultados ordenados con
  \passthrough{\lstinline!broom::tidy()!}.
\end{enumerate}

\section{Desarrollo}

\begin{lstlisting}[language=R]
library(tidyverse)
ventas  <- read_csv("datos/ventas.csv", show_col_types = FALSE)
tiendas <- read_csv("datos/tiendas.csv", show_col_types = FALSE)
ventas_det <- ventas |>
  left_join(tiendas, by = "tienda_id") |>
  drop_na(unidades) |>
  mutate(temporada = if_else(month(fecha) %in% 11:12, "Nov-Dic", "Resto"))
\end{lstlisting}

\subsection{Descriptivos}

\begin{lstlisting}[language=R]
ventas_det |>
  summarise(
    media   = mean(unidades),
    mediana = median(unidades),
    de      = sd(unidades),
    p90     = quantile(unidades, 0.9),
    .by = producto
  )
\end{lstlisting}

\subsection{Prueba t: ¿la temporada cambia las ventas?}

\begin{lstlisting}[language=R]
reloj <- ventas_det |> filter(producto == "Reloj")
t.test(unidades ~ temporada, data = reloj)       # Welch por defecto
wilcox.test(unidades ~ temporada, data = reloj)  # alternativa no paramétrica
\end{lstlisting}

\begin{quote}
\textbf{Cuidado:} las observaciones semanales de una misma tienda no son
del todo independientes. Aquí la prueba es ilustrativa; en un análisis
formal hay que considerar la autocorrelación.
\end{quote}

\subsection{Chi-cuadrada: ¿la mezcla de productos depende del canal?}

\begin{lstlisting}[language=R]
tabla <- xtabs(unidades ~ canal + producto, data = ventas_det)
tabla
chisq.test(tabla)
\end{lstlisting}

\subsection{Regresión lineal}

\begin{lstlisting}[language=R]
semanal <- ventas_det |>
  summarise(unidades = sum(unidades), .by = c(fecha, canal, temporada))

modelo <- lm(unidades ~ temporada + canal, data = semanal)
summary(modelo)
confint(modelo)

library(broom)
tidy(modelo)          # coeficientes como tibble
glance(modelo)        # R², AIC, etc.

nuevo <- tibble(temporada = "Nov-Dic", canal = "En línea")
predict(modelo, newdata = nuevo, interval = "prediction")
\end{lstlisting}

\subsection{Diagnóstico}

\begin{lstlisting}[language=R]
aug <- augment(modelo)
ggplot(aug, aes(.fitted, .resid)) +
  geom_point(alpha = 0.5) +
  geom_hline(yintercept = 0, linetype = "dashed") +
  labs(title = "Residuos vs. ajustados", x = "Ajustado", y = "Residuo")
\end{lstlisting}

\section{Actividades y ejercicios}

\begin{enumerate}
\def\labelenumi{\arabic{enumi}.}
\tightlist
\item
  Calcula la matriz de correlación entre
  \passthrough{\lstinline!precio!} y \passthrough{\lstinline!unidades!}
  a nivel registro. ¿Qué signo esperas y por qué?
\item
  Prueba si las unidades de \textbf{Audífonos} difieren entre las
  regiones \textbf{Norte} y \textbf{Sur}.
\item
  Ajusta \passthrough{\lstinline!lm(unidades \~ temporada * canal)!} y
  compara su R² ajustado con el modelo sin interacción.
\item
  Redacta en 3 líneas la interpretación del coeficiente
  \passthrough{\lstinline!temporadaResto!} para un público no técnico.
\end{enumerate}

\subsection{Soluciones}

\begin{lstlisting}[language=R]
# 1: negativo; los productos caros se venden en menos unidades
cor(ventas_det$precio, ventas_det$unidades)
cor(ventas_det$precio, ventas_det$unidades, method = "spearman")

# 2
audif <- ventas_det |> filter(producto == "Audífonos", region %in% c("Norte", "Sur"))
t.test(unidades ~ region, data = audif)

# 3
modelo_int <- lm(unidades ~ temporada * canal, data = semanal)
c(sin_interaccion = glance(modelo)$adj.r.squared,
  con_interaccion = glance(modelo_int)$adj.r.squared)
anova(modelo, modelo_int)

# 4 (ejemplo): "Fuera de noviembre-diciembre, cada canal vende en promedio
# alrededor de X unidades menos por semana que en temporada alta, manteniendo
# el canal constante; la diferencia es estadísticamente clara (p < 0.001)."
coef(modelo)["temporadaResto"]
\end{lstlisting}

\section{Evaluación (sumativa parcial)}

Informe breve (1 página) con: pregunta de negocio, prueba o modelo
elegido y \textbf{por qué}, resultado con intervalo de confianza,
verificación de supuestos y una limitación. Rúbrica: pertinencia del
método (30 \%), ejecución correcta (30 \%), interpretación (30 \%),
comunicación (10 \%).

\section{Referencias verificadas}

\begin{itemize}
\tightlist
\item
  \textbf{Video:} StatQuest with Josh Starmer. \emph{Linear Regression
  in R, Step-by-Step} {[}Video{]}. YouTube.
  \url{https://www.youtube.com/watch?v=u1cc1r_Y7M0} {[}EN{]}
  \emph{{[}fecha POR VERIFICAR{]}}
\item
  \textbf{Video:} StatQuest with Josh Starmer. \emph{Linear Regression
  and Linear Models} {[}Lista de reproducción{]}. YouTube.
  \url{https://www.youtube.com/playlist?list=PLblh5JKOoLUIzaEkCLIUxQFjPIlapw8nU}
  {[}EN{]}
\item
  \textbf{Web:} R Core Team. \emph{An Introduction to R} --- sección
  ``Statistical models in R''.
  \url{https://cran.r-project.org/doc/manuals/r-release/R-intro.html}
  {[}EN{]} {[}OA{]}
\item
  \textbf{Libro:} Wickham, H., Çetinkaya-Rundel, M., \& Grolemund, G.
  (2023). \emph{R for Data Science} (2.ª ed.). O'Reilly Media. ISBN
  978-1-4920-9740-2. Español: \url{https://davidrsch.github.io/r4ds-es/}
  {[}ES{]} {[}OA{]} --- caps. de análisis exploratorio como base para
  modelar.
\end{itemize}
